\chapter{Grundlagen und Stand der Technik}
\label{cha:grundlagen}

% Hier werden Begriffe definiert, die später gebraucht werden, und der aktuelle
% Forschungs-/Technikstand mit Quellen dargestellt. Nur das aufnehmen, was für
% Konzept, Umsetzung oder Evaluation tatsächlich benötigt wird.
% Richtwert: 20–30 % des Umfangs. Jeder Absatz mit Fremdwissen braucht ein Zitat.

\section{Begriffsdefinitionen}
\label{sec:begriffe}

% Abkürzungen mit \ac{…} einführen – beim ersten Vorkommen wird die Langform
% automatisch ausgeschrieben.
Eine \ac{API} beschreibt die Schnittstelle, über die Programme miteinander
kommunizieren~\cite{kopka2003guide}. Im Kontext dieser Arbeit wird unter \ldots{}
verstanden, dass \ldots

\section{Theoretische Grundlagen}
\label{sec:theorie}

% Modelle, Verfahren, Formeln. Gleichungen nummerieren und im Text referenzieren.
Die Laufzeit des Verfahrens lässt sich durch Gleichung~\eqref{eq:laufzeit}
abschätzen:
\begin{equation}
  T(n) = a \cdot n \log n + b
  \label{eq:laufzeit}
\end{equation}
Dabei bezeichnet $n$ die Eingabegröße und $a$, $b$ sind konstante Faktoren.

\section{Verwandte Arbeiten}
\label{sec:verwandte-arbeiten}

% Wer hat ähnliche Probleme wie gelöst? Was ist die Lücke, die diese Arbeit
% schließt? Ein Vergleich in Tabellenform ist hier oft hilfreich.
\textcite{knuth1984literate} beschreibt mit \emph{Literate Programming} einen Ansatz,
Quelltext und Dokumentation zu verbinden. Tabelle~\ref{tab:verwandte} stellt die
betrachteten Ansätze gegenüber.

\begin{table}[htbp]
  \centering
  \caption{Vergleich verwandter Ansätze}
  \label{tab:verwandte}
  \begin{tabular}{lccc}
    \toprule
    Ansatz & Lokal & Versionierung & Offen \\
    \midrule
    Webbasierter Editor & -- & \checkmark & -- \\
    Lokale IDE mit Plugin & \checkmark & (\checkmark) & \checkmark \\
    Diese Arbeit & \checkmark & \checkmark & \checkmark \\
    \bottomrule
  \end{tabular}
\end{table}

\section{Zusammenfassung}
% Kurz: Was folgt aus den Grundlagen für die eigene Arbeit?
Aus den betrachteten Arbeiten folgt, dass \ldots{} Diese Lücke adressiert das in
Kapitel~\ref{cha:umsetzung} vorgestellte Konzept.
